\documentclass[10pt,a4paper]{article}
\usepackage[top=0.85in, left=1.5in, right=1.5in, bottom=1.2in]{geometry}
\usepackage[utf8]{inputenc}
\usepackage[T1]{fontenc}
\usepackage{amsmath,amssymb,amsthm}
\usepackage{mathtools}
\usepackage{algorithm}
\usepackage{algpseudocode}
\usepackage{hyperref}
\usepackage{booktabs}
\usepackage{array}

\newtheorem{theorem}{Theorem}
\newtheorem{lemma}[theorem]{Lemma}
\newtheorem{proposition}[theorem]{Proposition}
\newtheorem{corollary}[theorem]{Corollary}
\newtheorem{remark}[theorem]{Remark}

\title{\textbf{Hybrid Quantum 4-Sieve}}
\author{
    MohamadAli Khajeian\\[0.5em]
    \normalsize School of Engineering Sciences, College of Engineering, University of Tehran\\
    \normalsize \texttt{khajeian@ut.ac.ir}
}
\date{\today}

\begin{document}
\maketitle

\begin{abstract}
We present an adaptation of the quantum random walk (QRW) 2-sieve to serve as the inner two-list subroutine inside a quantum 4-sieve. Inspired by the classical 4-sieve, we construct the quantum 4-sieve and use the QRW 2-sieve as its inner two-list subroutine. The adaptation preserves the Johnson graph, the second-level random product code, and the quantum data structure of the original algorithm, and modifies only the marked condition and the probability that a random cross-pair is a solution. We provide the complete algorithms, discuss the modifications required, and give a full complexity analysis for the running time, QRAM, quantum memory, and classical memory. We optimise the QRW parameters for the first two inner calls, and also analyse the third inner call in detail...
\end{abstract}

\section{Introduction}

The Shortest Vector Problem (SVP) is at the foundation of lattice-based cryptography. The best known heuristic quantum algorithm for SVP runs in time \(2^{0.2570d+o(d)}\) \cite{CL21}, obtained by replacing Grover's algorithm inside the 2-sieve by a quantum random walk with an additional layer of locality-sensitive filtering. In parallel, \(k\)-sieve algorithms have been developed to reduce memory requirements at the cost of increased running time \cite{CL23}. In particular, the quantum 4-sieve of \cite{CL23} uses nested Grover searches and achieves a time-memory trade-off with quantum-accessible classical memory of size \(2^{0.1724d}\) to \(2^{0.1813d}\).

In this work, we investigate whether the QRW 2-sieve of \cite{CL21} can be used as the inner two-list subroutine inside the 4-sieve framework of \cite{CL23}. This requires adapting the QRW 2-sieve to handle two distinct lists rather than one, an arbitrary configuration threshold \(C'_{12}(\alpha)\) rather than the fixed angle \(\theta_\alpha^*\), and possibly non-unit vectors in the second-level call.

We show that the core QRW structure remains unchanged: the Johnson graph, the second random product code \(\mathcal{C}_2\), the quantum data structure for the lists \(J^v(\vec{t})\), and the MNRS framework are all preserved. Only the marked condition and the probability \(p_{\mathrm{sol}}\) that a random cross-pair is a solution must be adjusted. We give the complete algorithms, discuss the modifications, and provide the full complexity analysis.

\section{Preliminaries and Notation}

We use the notation of \cite{CL21,CL23}. Let \(d\) be the lattice dimension and define the fundamental list-size parameter
\begin{equation}
N \coloneqq 2^{cd+o(d)}.
\end{equation}
This quantity \(N\) is exponential in \(d\) and serves as the reference scale for all exponents in the algorithm. A list of size \(N^{a}\) contains \(2^{a\cdot cd + o(d)}\) elements; throughout the paper we write \(N^{a+o(1)}\) to denote such a list, absorbing subexponential factors into the \(o(1)\).
For an angle \(\alpha \in [0,\pi/2]\), the volume ratio of a spherical cap of angle \(\alpha\) is
\begin{equation}
\mathcal{V}_d(\alpha) \coloneqq \mathrm{poly}(d)\sin^d(\alpha),
\end{equation}
which represents the fraction of the unit sphere covered by a cap of angular radius \(\alpha\). For two angles \(\alpha\) and \(\theta\), the volume ratio of a wedge formed by two caps of angle \(\alpha\) with centers separated by angle \(\theta\) is
\begin{equation}
\mathcal{W}_d(\alpha,\theta) \coloneqq \mathrm{poly}(d)\left(1-\frac{2\cos^2(\alpha)}{1+\cos(\theta)}\right)^{d/2},
\end{equation}
which is well-defined provided \(2\cos^2(\alpha) < 1+\cos(\theta)\), i.e., the two caps overlap.
The critical angle between residual vectors that characterises reducibility of two border points is
\begin{equation}
\theta_\alpha^* \coloneqq 2\arcsin\!\left(\frac{1}{2\sin(\alpha)}\right),
\end{equation}
as established in \cite[Proposition~4]{CL21}. Two residual vectors \(\vec{y}_0,\vec{y}_1\) on the border of a filter of angle \(\alpha\) are reducible if and only if \(\theta(\vec{y}_0,\vec{y}_1) \le \theta_\alpha^*\).

\paragraph{Outer 4-sieve parameters.} The outer 4-sieve uses a prefiltering angle \(\alpha_0 \in [\pi/3,\pi/2]\) and a configuration matrix \(C \in \mathbb{R}^{4\times 4}\). The reduction to residual vectors is denoted \(C'(\alpha_0)\), with entries \(C'_{ij}(\alpha_0)\) given by \cite[Lemma~1]{CL23}:
\begin{equation}
C'_{ij}(\alpha_0) = \frac{1}{\sin^2(\alpha_0)}\left(C_{ij} + \frac{\cos^2(\alpha_0)}{3}\right) \quad (i \neq j).
\end{equation}
The effective threshold for the second-level 2-sieve that builds \(L_{1234}\) is
\begin{equation}
C_{\mathrm{eff}} \coloneqq \frac{\sin^2(\alpha_0)}{2R^2} - 1, \qquad R^2 \coloneqq 2 + 2C_{12},
\end{equation}
as given in \cite[Lemma~6]{CL23}.

\paragraph{QRW parameters.} The adapted QRW 2-sieve uses four free parameters \(c, c_1, c_2, \rho\) satisfying \(0 \le c_2 \le c_1 \le c \le 1\) and \(0 \le \rho \le \rho_0\). From these, the following derived quantities are defined. The angle \(\alpha\) (not to be confused with \(\alpha_0\)) satisfies \(\mathcal{V}_d(\alpha) = N^{-(1-c)}\). The angle \(\beta\) of the second random product code satisfies \(\mathcal{V}_d(\beta) = N^{c_2-c_1}\). The target angle is \(\theta' \coloneqq \arccos(C')\). The collision parameter \(\rho_0\) is defined by \(N^{\rho_0} \coloneqq \mathcal{V}_d(\beta)/\mathcal{W}_d(\beta,\theta')\). The solution probability is
\begin{equation}
p_{\mathrm{sol}} \coloneqq \frac{1}{2}\bigl(1-(C')^2\bigr)^{d/2},
\label{eq:psol}
\end{equation}
where the factor \(1/2\) accounts for the bipartite nature of the problem. The expected number of solutions per filter is \(N^\zeta \coloneqq N^{2c}\cdot p_{\mathrm{sol}}\).

The QRW costs are the setup cost \(\mathcal{S} \coloneqq N^{c_1+\rho}\), the update cost \(\mathcal{U} \coloneqq N^{\max\{\rho,(\rho+c_2)/2\}}\), and the checking cost \(\mathcal{C} \coloneqq 1\). The spectral gap of the Johnson graph is \(\delta \coloneqq N^{-c_1}\). The fraction of marked vertices is
\begin{equation}
\epsilon \coloneqq N^{2c_1}\cdot p_{\mathrm{sol}}\cdot N^{\rho-\rho_0}.
\end{equation}

\section{Modified QRW 2-Sieve}

\subsection{Preserved components}

The adaptation preserves the following components of the QRW 2-sieve of \cite{CL21} without structural modification. The Johnson graph \(J(N, N^{c_1})\) on the combined list, with spectral gap \(\delta = N^{-c_1}\) and the MNRS quantum walk framework, is retained exactly. The second random product code \(\mathcal{C}_2\) with angle \(\beta\) satisfying \(\mathcal{V}_d(\beta) = N^{c_2-c_1}\) and size \(|\mathcal{C}_2| = N^{\rho+c_1-c_2}\) is constructed identically. The quantum data structure for the lists \(J^v(\vec{t}) = f_\beta(\vec{t})\cap L^v\), supporting insertions and deletions in quantum superposition, is used as in the original algorithm. The setup, update, and check cost formulas are formally unchanged.

\subsection{Required modifications}

Three modifications are required to adapt the QRW 2-sieve to the bipartite setting. First, we form the combined list \(L = R_A \cup R_B\) and tag each element by a bit indicating its origin. The Johnson graph is built on \(L\) with \(|L| = 2N_R\). This tagging does not affect the graph structure. Second, the marked condition is redefined: a vertex \(v\) is marked if there exist \(\vec{t} \in \mathcal{C}_2\) and distinct \(\vec{y}_i,\vec{y}_j \in \tilde{J}^v(\vec{t})\) such that \(\vec{y}_i\) and \(\vec{y}_j\) have different tags and \(\langle \vec{y}_i|\vec{y}_j\rangle \le C'\). Third, the solution probability acquires the factor \(1/2\) as given in Equation~\eqref{eq:psol}.

\subsection{Adapted algorithm}

\begin{algorithm}[H]
\caption{Adapted QRW 2-Sieve for bipartite configuration problems}
\label{alg:qrw}
\begin{algorithmic}[1]
\Require two residual lists \(R_A, R_B\) of size \(N_R\) each; threshold \(C'\); QRW parameters \(c, c_1, c_2, \rho\).
\Ensure all pairs \((\vec{y}_a, \vec{y}_b) \in R_A \times R_B\) with \(\langle \vec{y}_a|\vec{y}_b\rangle \le C'\).
\State Form \(L \coloneqq R_A \cup R_B\) with tags.
\State Compute \(\theta' \coloneqq \arccos(C')\) and \(p_{\mathrm{sol}} \coloneqq \frac{1}{2}(1-C'^2)^{d/2}\).
\State Sample random product code \(\mathcal{C}_2\) with \(|\mathcal{C}_2| = N^{\rho+c_1-c_2}\) and angle \(\beta\) satisfying \(\mathcal{V}_d(\beta) = N^{c_2-c_1}\).
\State Build Johnson graph \(J(N, N^{c_1})\) on \(L\).
\For{each vertex \(v\)}
    \State Compute \(J^v(\vec{t}) = f_\beta(\vec{t}) \cap L^v\) for each \(\vec{t}\in\mathcal{C}_2\).
    \State Truncate to \(\tilde{J}^v(\vec{t})\) of size \(\le 2N^{c_2}\).
    \State Set marked bit \(M(v)\) according to the marked condition above.
\EndFor
\While{total solutions found \(< N^\zeta/2\)}
    \State Pick a new random product code \(\mathcal{C}_2\).
    \While{solutions found with this \(\mathcal{C}_2\) \(< N^{\zeta+\rho-\rho_0}/2\)}
        \State Run the MNRS quantum walk to find one marked vertex.
    \EndWhile
\EndWhile
\State \Return all collected pairs.
\end{algorithmic}
\end{algorithm}

\section{Hybrid Quantum 4-sieve}

\begin{algorithm}[H]
\caption{Hybrid Quantum 4-sieve}
\label{alg:main}
\begin{algorithmic}[1]
\Require list \(L\) of \(N' = N^{1+o(1)}\) lattice vectors of norm \(1\); reducing factor \(\gamma < 1\); prefiltering angle \(\alpha_0\); configuration \(C\in\mathbb{R}^{4\times 4}\); QRW parameters.
\Ensure list \(L'\) of \(N'\) lattice vectors of norm \(\le \gamma\).
\State \(L' \coloneqq \emptyset\).
\While{\(|L'| \le N'\)}
    \State Sample a 4-RPC \(\mathfrak{C}\) with codewords \((\mathbf{A}_1^j,\mathbf{A}_2^j,\mathbf{A}_3^j,\mathbf{A}_4^j)\) such that \(\sum_i \mathbf{A}_i^j = \vec{0}\) and \(\langle \mathbf{A}_i^j|\mathbf{A}_{i'}^j\rangle = -1/3\) for \(i\neq i'\).
    \State Partition \(L\) into tuple-filters: for each \(\vec{x}\in L\), find its nearest codeword \(\mathbf{A}_i^j\), add \(\vec{x}\) to \(L_i^j\), compute residual \(\vec{y} = (\vec{x}-\cos\alpha_0\,\mathbf{A}_i^j)/\sin\alpha_0\), and add \(\vec{y}\) to \(R_i^j\).
    \For{each tuple-filter \(j\)}
        \State \(L_{12}^j \leftarrow \textsc{QRW-2-Sieve}(R_1^j, R_2^j, C'_{12}(\alpha_0))\).
        \State \(L_{34}^j \leftarrow \textsc{QRW-2-Sieve}(R_3^j, R_4^j, C'_{34}(\alpha_0))\).
        \State \(L_{1234}^j \leftarrow \textsc{Grover}(\tilde{L}_{12}^j, \tilde{L}_{34}^j, C_{\mathrm{eff}})\).
        \State \(L' \coloneqq L' \cup L_{1234}^j\).
    \EndFor
\EndWhile
\State \Return \(L'\).
\end{algorithmic}
\end{algorithm}

\section{Complexity Analysis}

We now give the full complexity of the adapted QRW 2-sieve and of the outer 4-sieve. All calculations are performed in the QRAM model and expressed in the exponent notation \(N^{cd+o(1)}\).

\subsection{Complexity of the adapted QRW 2-sieve}

The setup cost, update cost, checking cost, and spectral gap are formally unchanged from the original algorithm:
\begin{align}
\mathcal{S} &= N^{c_1+\rho+o(1)}, \\
\mathcal{U} &= N^{\max\{\rho,(\rho+c_2)/2\}+o(1)}, \\
\mathcal{C} &= 1, \\
\delta &= N^{-c_1}.
\end{align}
The fraction of marked vertices is modified by the factor \(1/2\) in \(p_{\mathrm{sol}}\):
\begin{equation}
\epsilon = \frac{1}{2}\,N^{2c_1}\bigl(1-(C')^2\bigr)^{d/2}\,N^{\rho-\rho_0},
\end{equation}
where \(\rho_0\) satisfies \(N^{\rho_0} = \mathcal{V}_d(\beta)/\mathcal{W}_d(\beta,\theta')\). The expected number of solutions per filter is
\begin{equation}
N^\zeta = \frac{1}{2}\,N^{2c}\bigl(1-(C')^2\bigr)^{d/2}.
\end{equation}
The time to find one marked vertex is
\begin{equation}
T_1 = \mathcal{S} + \frac{1}{\sqrt{\epsilon}}\left(\frac{1}{\sqrt{\delta}}\,\mathcal{U} + \mathcal{C}\right).
\label{eq:T1}
\end{equation}
The time for the \textsc{FindAllSolutions} subroutine is
\begin{equation}
\mathrm{FAS}_1 = \max\{N^\zeta, 1\}\cdot T_1.
\label{eq:FAS1}
\end{equation}
The total time of one QRW 2-sieve call is
\begin{equation}
T_{\mathrm{QRW}} = \mathrm{NbRep}_{\mathrm{QRW}}\cdot\bigl(N + N^{1-c}\cdot \mathrm{FAS}_1\bigr),
\label{eq:TQRW}
\end{equation}
where
\begin{equation}
\mathrm{NbRep}_{\mathrm{QRW}} = \max\bigl\{1,\, N^{c-\zeta+o(1)}\bigr\}.
\label{eq:NbRepQRW}
\end{equation}
The derived parameters are:
\begin{align}
\log_2 \sin(\alpha) &= -(1-c)\cdot n, \\
\log_2 \sin(\beta) &= (c_2-c_1)\cdot n, \\
\theta' &= \arccos(C'), \\
\rho_0 &= \frac{\log_2 \mathcal{V}_d(\beta) - \log_2 \mathcal{W}_d(\beta,\theta')}{n\,d}, \\
\zeta &= 2c + \frac{\log_2 p_{\mathrm{sol}}}{n\,d}.
\end{align}

\subsection{QRW 2-sieve on the standard 2-sieve problem}

We first verify the framework by reproducing the result of \cite[Theorem~1]{CL21} for the standard 2-sieve problem. In this case, the list size is \(N\) (i.e., \(n = 0.2075\)), the threshold is \(\theta_\alpha^*\) with \(p_{\mathrm{sol}} = \mathcal{V}_d(\theta_\alpha^*)\), and the optimal parameters are
\begin{equation}
c \approx 0.3696, \quad c_1 \approx 0.2384, \quad c_2 = 0, \quad \rho \to 0.
\end{equation}
With these values:
\begin{align}
\log_2 \sin(\alpha) &= -0.1308, & \alpha &\approx 1.1514 \text{ rad}, \\
\log_2 \sin(\beta) &= -0.0495, & \beta &\approx 1.313 \text{ rad}, \\
\theta_\alpha^* &\approx 1.158 \text{ rad}, & \rho_0 &\approx 0.106, & \zeta &\approx 0.132.
\end{align}
The QRW costs are:
\begin{align}
\mathcal{S} &= N^{0.2384}, & \mathcal{U} &= N^{o(1)}, & \delta &= N^{-0.2384}, & \epsilon &= N^{-0.2364}.
\end{align}
Then \(T_1 \approx N^{0.2384}\), \(\mathrm{FAS}_1 = N^{0.3704}\), and
\begin{equation}
T_{\mathrm{QRW}} \approx N^{1.2384} = 2^{0.2570d+o(d)},
\end{equation}
which matches \cite[Theorem~1]{CL21}.

\subsection{QRW 2-sieve for the first two inner calls}
\label{sec_QRW_first2}

We now apply the framework to the first two inner calls of the 4-sieve at the minimal-memory configuration of \cite[Proposition~17]{CL23}: \(|L| = 2^{0.1724d}\), \(\alpha_0 \approx 1.3131\) rad, and balanced configuration \(C_{ij} = -1/4\) for \(i \neq j\).
The residual list size is \(|L_i| = 2^{0.1240d}\). The threshold is
\begin{equation}
C'_{12}(\alpha_0) = \frac{-1/4 + \cos^2(1.3131)/3}{\sin^2(1.3131)} \approx -0.244.
\end{equation}
The solution probability is
\begin{equation}
p_{\mathrm{sol}} = \frac{1}{2}(0.9405)^{d/2} = \frac{1}{2}\cdot 2^{-0.0443d}.
\end{equation}
For the combined list \(N_1 = 2|L_i|\) with \(n_1 = 0.1240\), we have \(\log_{N_1} p_{\mathrm{sol}} \approx -0.360\).
We optimise the QRW parameters. The optimal choice in the plateau region is
\begin{equation}
c = 0.3, \quad c_1 = 0.13, \quad c_2 = 0, \quad \rho = 0.
\end{equation}
With these values, we compute the collision parameter carefully. We have
\begin{align}
\log_2 \sin(\beta) &= (c_2-c_1)\cdot n_1 = -0.13 \cdot 0.1229 = -0.0160, \\
\cos^2(\beta) &\approx 0.0219, \\
1 + \cos(\theta') &= 1 - 0.244 = 0.756, \\
\log_2 \mathcal{W}_d(\beta,\theta')/d &= \tfrac{1}{2}\log_2\!\Bigl(1-\tfrac{2\cos^2(\beta)}{1+\cos(\theta')}\Bigr) = -0.0431, \\
\rho_0 &= \frac{-0.0160 - (-0.0431)}{0.1229} \approx 0.2205.
\end{align}
The marked fraction is therefore
\begin{equation}
\epsilon = N_1^{2c_1 + \zeta - 2c + \rho - \rho_0} = N_1^{0.26 - 0.360 - 0.2205} = N_1^{-0.3205}.
\end{equation}
The remaining quantities are
\begin{align}
\zeta &= 2(0.3) - 0.360 = 0.240, \\
T_1 &= N_1^{0.13} + N_1^{0.1603}\bigl(N_1^{0.065} + 1\bigr) \approx N_1^{0.2253}, \\
\mathrm{FAS}_1 &= N_1^{0.240} \cdot N_1^{0.2253} = N_1^{0.4653}, \\
\mathrm{NbRep}_{\mathrm{QRW}} &= N_1^{0.3 - 0.240} = N_1^{0.060}, \\
T_{\mathrm{QRW}} &= N_1^{0.060}\bigl(N_1 + N_1^{0.7}\cdot N_1^{0.4653}\bigr) \approx N_1^{1.2253}.
\end{align}
Converting \(T_{12} = T_{34} = 2^{0.1506d+o(d)} \approx 2^{0.151d}\).


\subsection{Complexity of the Third Inner Call}
\label{sec:third-call}

After the first two inner calls, we have obtained the lists
\(L_{12}\) and \(L_{34}\) of pairs of residual vectors that satisfy the
first‑level configuration constraints \(C'_{12}(\alpha_0)\) and
\(C'_{34}(\alpha_0)\), respectively. The third inner call consists in
finding all pairs
\[
(\mathbf z_{12},\mathbf z_{34})\in L_{12}\times L_{34}
\]
such that
\begin{equation}
\langle \mathbf z_{12}|\mathbf z_{34}\rangle \le C_{\mathrm{eff}},
\qquad
C_{\mathrm{eff}} \coloneqq \frac{\sin^2(\alpha_0)}{2R^2}-1,
\qquad
R^2 \coloneqq 2+2C_{12},
\label{eq:ceff}
\end{equation}
where \(\alpha_0\) is the outer prefiltering angle and \(C_{12}\) is the
outer configuration entry. This condition is exactly the one given in
\cite[Lemma~6]{CL23} and guarantees that the four residual vectors form
a reducing tuple.

\paragraph{Probability of a solution.}
For two random vectors \(\mathbf z_{12}\) and \(\mathbf z_{34}\) of norm
\(R\), the probability that their inner product is at most
\(C_{\mathrm{eff}}\) is
\begin{equation}
p_{\mathrm{sol}} = \bigl(1-C_{\mathrm{eff}}^2\bigr)^{d/2}.
\label{eq:psol-third}
\end{equation}
There is no factor \(1/2\) here, because the two lists \(L_{12}\) and
\(L_{34}\) are distinct and the search is over ordered pairs
\((\mathbf z_{12},\mathbf z_{34})\).

The expected number of solutions per tuple‑filter is therefore
\begin{equation}
|\mathrm{Sol}| = |L_{12}|^2\,p_{\mathrm{sol}}
= |L_{12}|^2\bigl(1-C_{\mathrm{eff}}^2\bigr)^{d/2}.
\label{eq:sol-third}
\end{equation}

\paragraph{Grover search.}
We use Grover's algorithm to find all solutions in
\(L_{12}\times L_{34}\). The search space has size
\(|L_{12}|^2\). The number of Grover iterations to find one solution is
\(O\bigl(1/\sqrt{p_{\mathrm{sol}}}\bigr)\), and to find all
\(|\mathrm{Sol}|\) solutions we use amplitude amplification, which
requires
\begin{equation}
T_{1234}
= |L_{12}|^2\bigl(1-C_{\mathrm{eff}}^2\bigr)^{d/4}
= |L_{12}|^2\sqrt{p_{\mathrm{sol}}}
\label{eq:T1234}
\end{equation}
elementary operations.

\paragraph{Numerical values for the minimal‑memory configuration.}
We take the minimal‑memory parameters of \cite[Proposition~17]{CL23}:
\[
|L| = 2^{0.1724d},\qquad
\alpha_0 \approx 1.3131\ \text{rad},\qquad
C_{ij}=-\frac14\ (i\neq j).
\]
Then
\[
\sin^2(\alpha_0) \approx 0.9349,\qquad
R^2 = 1.5,\qquad
C_{\mathrm{eff}} \approx -0.6884.
\]
The first‑level threshold is
\[
C'_{12}(\alpha_0)
= \frac{-\frac14+\frac{\cos^2(\alpha_0)}{3}}{\sin^2(\alpha_0)}
\approx -0.2443.
\]
The residual list size is
\[
|R_1| = |L|\,\mathcal V_d(\alpha_0)
= 2^{0.1724d}\cdot 2^{-0.0486d}
= 2^{0.1238d},
\]
so that
\[
|L_{12}|
= |R_1|^2\bigl(1-C'_{12}(\alpha_0)^2\bigr)^{d/2}
= 2^{0.2033d}.
\]
The solution probability is
\[
p_{\mathrm{sol}}
= \bigl(1-C_{\mathrm{eff}}^2\bigr)^{d/2}
= 2^{-0.4635d},
\]
hence
\[
T_{1234}
= 2^{0.2033d\cdot 2}\cdot 2^{-0.4635d/2}
= 2^{0.2033d+o(d)}.
\]
From Section~\ref{sec_QRW_first2}, the first two inner calls
cost
\[
T_{12} = T_{34} \approx 2^{0.1517d+o(d)}.
\]
Thus the total time per tuple‑filter is dominated by the third call:
\begin{equation}
T_{\text{filter}}
= T_{12}+T_{34}+T_{1234}
\approx 2^{0.2033d+o(d)}.
\label{eq:Tfilter}
\end{equation}

Finally, we multiply by the number of tuple‑filters and the number of
repetitions of the outer loop. For the minimal‑memory configuration,
\cite[Proposition~17]{CL23} gives
\[
\mathrm{NbFilters} = 2^{0.0486d},\qquad
\mathrm{NbRep} = 2^{0.1069d}.
\]
Therefore the total running time of the hybrid quantum 4‑sieve is
\begin{equation}
T_{\text{hybrid}}
= \mathrm{NbRep}\cdot \mathrm{NbFilters}\cdot T_{\text{filter}}
= 2^{0.1069d}\cdot 2^{0.0486d}\cdot 2^{0.2033d}
= 2^{0.3588d+o(d)}.
\label{eq:Thybrid}
\end{equation}
We round this to
\[
\boxed{T_{\text{hybrid}} \approx 2^{0.359d+o(d)}.}
\]


\section{Comparison}
\label{sec:comparison}

We now compare the hybrid quantum 4‑sieve with the classical 4‑sieve
of \cite{CL23} at the same classical memory. The relevant data for the
minimal‑memory regime is summarized in Table~\ref{tab:comparison}.

\begin{table}[h]
\centering
\begin{tabular}{lccc}
\toprule
Algorithm & Classical memory & Running time \\
\midrule
Classical 4‑sieve \cite{CL23} & \(2^{0.1724d}\) & \(2^{0.418d}\) \\
Hybrid quantum 4‑sieve (ours) & \(2^{0.1724d}\) & \(2^{0.359d}\) \\
Grover 4‑sieve \cite{CL23} & \(2^{0.1724d}\) & \(2^{0.3276d}\) \\
\bottomrule
\end{tabular}
\caption{Comparison at minimal classical memory.}
\label{tab:comparison}
\end{table}

At the minimal classical memory \(M=2^{0.1724d}\), the classical
4‑sieve runs in time \(2^{0.418d+o(d)}\), whereas our hybrid algorithm
runs in time \(2^{0.359d+o(d)}\). The improvement factor is
\[
2^{0.418d-0.359d} = 2^{0.059d},
\]
which is an exponential speedup. For comparison, the pure quantum
4‑sieve of \cite{CL23} achieves \(2^{0.3276d+o(d)}\) at the same
memory, so our hybrid is only slightly slower (\(2^{0.0312d}\)) than
the best known quantum 4‑sieve.

The advantage of the hybrid is most pronounced in the low‑memory
regime. As the allowed memory increases, the classical 4‑sieve
improves rapidly: at \(M=2^{0.1887d}\) it runs in \(2^{0.324d}\), and at
\(M=2^{0.2075d}\) it reaches the classical 2‑sieve time
\(2^{0.2925d}\). Our hybrid, optimized for the minimal‑memory point,
does not immediately benefit from additional memory without a
re‑optimization of the QRW parameters. A full time‑memory trade‑off
curve for the hybrid would require running the parameter optimization
for several values of \(\alpha_0\) and \(C\), which we leave for future
work.

In summary, the hybrid quantum 4‑sieve demonstrates that quantum
random walk techniques can be successfully integrated into the
\(k\)-sieve framework. While it does not surpass the best pure quantum
4‑sieve, it offers a significant quantum advantage over the classical
4‑sieve in the regime where memory is severely constrained.


\begin{thebibliography}{9}
\bibitem{CL21} A.~Chailloux and J.~Loyer. Lattice sieving via quantum random walks. In \emph{ASIACRYPT 2021}, 2021.
\bibitem{CL23} A.~Chailloux and J.~Loyer. Classical and quantum 3 and 4-sieves to solve SVP with low memory. In \emph{PQCrypto 2023}, pp.~225--255, 2023.
\end{thebibliography}

\end{document}